\documentclass[11pt,a4paper]{article}

\usepackage[utf8]{inputenc}
\usepackage[T1]{fontenc}
\usepackage[brazil]{babel}
\usepackage{geometry}
\usepackage{hyperref}
\usepackage{enumitem}
\usepackage{booktabs}
\usepackage{xcolor}

\geometry{margin=2.5cm}
\hypersetup{colorlinks=true, linkcolor=blue!60!black, urlcolor=blue!60!black}

\title{\textbf{Atividade prática — Reinforcement Learning aplicado a combate BVR}\\[0.4em]
\large Disciplina TE-276}
\author{Plataforma online BVR Lab}
\date{Junho de 2026}

\begin{document}
\maketitle

\section*{Informações da disciplina}
\begin{tabular}{@{}ll@{}}
  \toprule
  \textbf{Disciplina} & TE-276 \\
  \textbf{Atividade} & Projeto de função de recompensa para missão ar-ar BVR \\
  \textbf{Plataforma} & \url{https://bvrlabonline.conceptio.com.br} \\
  \textbf{Código de turma} & Fornecido pelo professor (padrão: \texttt{BVR2026}) \\
  \bottomrule
\end{tabular}

\section{Objetivo}
Você \textbf{não} vai programar o algoritmo de aprendizado (PPO) nem alterar a rede neural.
Sua tarefa é atuar como \textbf{designer da missão}: definir pesos de recompensas e penalidades
para que um agente fixo aprenda a completar uma missão \emph{Beyond Visual Range} (BVR):
\textbf{alcançar a zona de objetivo vivo}, enfrentando oponentes inteligentes.

\section{O que está fixo vs.\ o que você controla}
\begin{tabular}{@{}p{0.45\textwidth}p{0.45\textwidth}@{}}
  \toprule
  \textbf{Você pode alterar} & \textbf{Bloqueado (não editar)} \\
  \midrule
  Pesos em recompensas e penalidades & Arquitetura da rede neural \\
  Seleção de oponentes para treino & Algoritmo PPO e hiperparâmetros \\
  Número de passos (dentro da cota) & Física do simulador \\
  & Oponentes oficiais da competição (B1--B10 + arquetipos) \\
  \bottomrule
\end{tabular}

\section{Pontuação oficial}
Após cada treino, o sistema avalia seu modelo contra \textbf{todos} os oponentes bloqueados.
A nota da competição usa:
\[
  \text{score} = 0{,}6 \times \text{missão} + 0{,}25 \times \text{abate} + 0{,}15 \times \text{eficiência de míssil}
\]
\begin{itemize}[nosep]
  \item \textbf{Missão}: chegar à zona verde \textbf{vivo} (objetivo principal).
  \item \textbf{Abate}: derrubar o inimigo ajuda a sobreviver, mas não substitui a missão.
  \item \textbf{Eficiência de míssil}: relação entre abates e mísseis disparados.
\end{itemize}

\section{Passo a passo na plataforma}

\subsection*{Passo 1 — Acesso e cadastro}
\begin{enumerate}[label=\textbf{\arabic*.}]
  \item Abra \url{https://bvrlabonline.conceptio.com.br}.
  \item Clique em \textbf{Register} e crie sua conta (nome, senha, código da turma).
  \item Faça \textbf{Login}. Observe sua \textbf{cota} no topo (número limitado de treinos por janela de tempo).
\end{enumerate}

\subsection*{Passo 2 — Configurar recompensas e oponentes (\emph{Train})}
No painel esquerdo da aba \textbf{Train}:
\begin{enumerate}[label=\textbf{\arabic*.},start=4]
  \item Ajuste os pesos de recompensa/penalidade (valor \textbf{0} desliga o termo).
  \item No mapa de oponentes, selecione quem usar no treino (B1--B10 e arquetipos). Prioridade 0 remove o oponente.
  \item Clique em \textbf{Queue training run}.
\end{enumerate}

\textbf{Dica:} comece com os valores padrão, rode uma vez, depois mude \textbf{um} termo por vez.

\subsection*{Passo 3 — Acompanhar o treino (\emph{Train}, painel direito)}
Durante a execução, observe:
\begin{itemize}[nosep]
  \item \textbf{progress} — avanço do treino;
  \item \textbf{avg score} — desempenho nas avaliações de monitoramento;
  \item \textbf{train reward} — recompensa média durante o aprendizado;
  \item \textbf{Live eval} — um card por oponente avaliado; clique em \textbf{▶ Play} no card para ver o replay (um por vez);
  \item \textbf{Abas Epoch} — clique em \emph{Epoch 1, 2, \ldots} para rever avaliações anteriores; use \emph{live} para a avaliação em andamento;
  \item \textbf{Learning curve} — linha verde subindo indica que sua função de recompensa está funcionando.
\end{itemize}

\subsection*{Passo 4 — Analisar o run concluído}
Quando o status for \textbf{done}, o painel \textbf{Run complete} e o relatório em \textbf{My runs} $\rightarrow$ \textbf{view report} mostram automaticamente:
\begin{enumerate}[label=\textbf{\arabic*.},start=7]
  \item \textbf{Parâmetros usados} — oponentes, passos e pesos de recompensa.
  \item \textbf{Curva de aprendizado} e \textbf{relatório de análise} — score, tabela por oponente e gráficos (contribuição das recompensas, eficiência de míssil, etc.).
  \item \textbf{Mapa de perfil dos agentes} — posição ofensiva/defensiva dos oponentes e estimativa de onde seu agente aprendido se situa (★).
  \item \textbf{Eval replay} — escolha \textbf{um oponente por vez} no menu e clique em \textbf{Watch replay}. Se os episódios de avaliação tiverem resultados diferentes, o replay mostra o \textbf{pior caso} (útil para diagnosticar falhas).
\end{enumerate}

\subsection*{Passo 5 — Enviar o melhor modelo (\emph{My runs})}
\begin{enumerate}[label=\textbf{\arabic*.},start=11]
  \item Vá à aba \textbf{My runs}.
  \item Compare \textbf{score}, \textbf{mission}, \textbf{kill} e \textbf{reward}.
  \item Clique em \textbf{submit} no seu melhor run (\textbf{apenas um} conta no ranking).
  \item Use \textbf{view report} para abrir o relatório completo de qualquer execução.
\end{enumerate}

\subsection*{Passo 6 — Preencher o relatório (\emph{Report})}
\begin{enumerate}[label=\textbf{\arabic*.},start=15]
  \item Abra a aba \textbf{Report}.
  \item Responda às \textbf{2 perguntas} (veja seção~\ref{sec:relatorio}).
  \item Clique em \textbf{Salvar relatório} antes do prazo.
\end{enumerate}

\subsection*{Passo 7 — Ranking (\emph{Leaderboard})}
\begin{enumerate}[label=\textbf{\arabic*.},start=18]
  \item Confira sua posição na aba \textbf{Leaderboard}.
  \item O ranking usa o score oficial contra o conjunto bloqueado de oponentes.
\end{enumerate}

\section{Termos de recompensa (referência rápida)}
\begin{tabular}{@{}p{0.28\textwidth}p{0.62\textwidth}@{}}
  \toprule
  \textbf{Termo} & \textbf{Significado} \\
  \midrule
  \texttt{mission\_completed} & Bônus grande por completar a missão vivo \\
  \texttt{hit\_enemy} & Bônus por abater o inimigo \\
  \texttt{was\_hit} & Penalidade ao ser abatido (mantenha negativo) \\
  \texttt{fire\_missile} / \texttt{miss\_missile} & Custo de disparo / tiro errado \\
  \texttt{mission\_shaping} & Recompensa por progresso em direção ao objetivo \\
  \texttt{maintain\_track} / \texttt{lost\_track} & Manter / perder radar no alvo \\
  \texttt{closing\_bonus} / \texttt{wez\_advantage} & Aproximação / vantagem de alcance \\
  \bottomrule
\end{tabular}

\section{Entregáveis}
\begin{enumerate}[label=\textbf{E\arabic*.}]
  \item \textbf{Modelo submetido} — melhor run enviado via \emph{My runs} $\rightarrow$ \emph{submit}.
  \item \textbf{Relatório} — formulário da aba \emph{Report} (2 respostas objetivas).
  \item \textbf{Participação no ranking} — score visível na aba \emph{Leaderboard}.
\end{enumerate}

\section{Relatório simplificado}\label{sec:relatorio}
Preencha com base nos gráficos, no mapa de perfil, nos replays e no run submetido. Seja \textbf{objetivo}.

\begin{enumerate}[label=\textbf{R\arabic*.}]
  \item \textbf{Estratégia de Feedbacks} — Como você ajustou os pesos de recompensa e penalidade entre cada treino? Que comportamento esperava e o que obteve?
  \item \textbf{Análise dos Resultados Finais} — Analise o comportamento do run submetido (score, curva, mapa de perfil, replays, oponentes) e sugira o que tornaria o agente mais robusto.
\end{enumerate}

\section{Erros comuns}
\begin{itemize}[nosep]
  \item Priorizar abate em detrimento da missão $\rightarrow$ agente luta, mas não completa o objetivo.
  \item Zerar todos os termos de \emph{shaping} $\rightarrow$ aprendizado muito lento no início.
  \item Ignorar replays $\rightarrow$ perde a causa das derrotas (rota errada, disparo precoce, etc.).
  \item Esquecer de \textbf{submit} e \textbf{salvar} o relatório.
\end{itemize}

\section{Suporte}
\begin{tabular}{@{}ll@{}}
  \toprule
  Problema & Ação \\
  \midrule
  Cota esgotada & Aguardar a janela de tempo ou falar com o professor \\
  Página desatualizada & Atualização forçada: Ctrl+F5 \\
  Replay no Live eval não anima & Clique em \textbf{▶ Play} no card do oponente desejado \\
  Replay no relatório não aparece & Aguarde \emph{report ready}; só runs concluídos após a atualização têm replays gravados \\
  \bottomrule
\end{tabular}

\vfill
\begin{center}
  \small \textbf{Disciplina TE-276} — BVR Reinforcement Learning Lab\\
  Dashboard local (opcional): \url{https://bvrlab.conceptio.com.br}
\end{center}

\end{document}
